% Exact geometry: P = [0,1]^2 intersect {z_1 + z_2 <= 3/2}.
% Canvas: 112 x 72 mm; labels inherit the volume's WenKai figure style.
\begin{tikzpicture}[x=1mm,y=1mm,font=\figurefont\footnotesize,
  text=FigureInk,draw=FigureStroke,line width=.8pt]
  \path[use as bounding box] (0,0) rectangle (112,72);
  \begin{scope}[shift={(12,12)},x=50mm,y=50mm]
    \fill[FigureYellowFill] (0,0)--(1,0)--(1,.5)--(.5,1)--(0,1)--cycle;
    \draw[FigureGrid,dashed,line width=.6pt] (1,0)--(1,1)--(0,1);
    \draw (0,0)--(1,0)--(1,.5)--(.5,1)--(0,1)--cycle;
    \draw[FigureRed,line width=1.3pt] (.5,1)--(1,.5);
    \draw[-{Stealth[length=1.8mm,width=1.2mm]}] (-.06,0)--(1.15,0);
    \draw[-{Stealth[length=1.8mm,width=1.2mm]}] (0,-.06)--(0,1.12);
    \foreach \x in {.5,1} {
      \draw (\x,-.018)--(\x,.018);
      \draw (-.018,\x)--(.018,\x);
    }
    \node[below left] at (0,0) {$0$};
    \node[below=3pt] at (.5,0) {$1/2$};
    \node[below=3pt] at (1,0) {$1$};
    \node[left=3pt] at (0,.5) {$1/2$};
    \node[left=3pt] at (0,1) {$1$};
    \node[right] at (1.15,0) {$z_1$};
    \node[above] at (0,1.12) {$z_2$};
    \node at (.28,.76) {$P$};
    \foreach \p in {(0,0),(1,0),(0,1)} {
      \filldraw[draw=FigureStroke,fill=FigureBlue] \p circle (1.1mm);
    }
    \fill[FigureMuted] (.23,.23) rectangle (.27,.27);
    \node[above right=2pt] at (.25,.25) {$G=1/2$};
    \draw[FigureRed,line width=1.1pt] (1,.5) circle (1.3mm);
    \node[right=5pt] at (1,.5) {$G=3/2$};
    \draw[FigureStroke,line width=1pt] (.976,.976)--(1.024,1.024);
    \draw[FigureStroke,line width=1pt] (.976,1.024)--(1.024,.976);
  \end{scope}
  \node[anchor=west] at (70,61) {$z_1+z_2=3/2$};
  \draw[FigureMuted,line width=.6pt] (68,59)--(51,51);
  \filldraw[draw=FigureStroke,fill=FigureBlue] (82,27) circle (1.1mm);
  \node[anchor=west] at (85,27) {原可行点};
  \fill[FigureMuted] (81,21) rectangle (83,23);
  \node[anchor=west] at (85,22) {松弛候选};
  \draw[FigureRed,line width=1.1pt] (82,17) circle (1.3mm);
  \node[anchor=west] at (85,17) {松弛最优点};
  \draw[FigureStroke,line width=1pt] (81,11)--(83,13);
  \draw[FigureStroke,line width=1pt] (81,13)--(83,11);
  \node[anchor=west] at (85,12) {不可行点};
\end{tikzpicture}%
